\section{Background}
\label{sec:background}

\subsection{CVE and CWE}

The CVE program, operated by MITRE on behalf of the U.S.\ government, assigns
a unique identifier of the form \texttt{CVE-YYYY-NNNNN} to each publicly
disclosed vulnerability, together with a short description and a list of
references~\cite{cve}. The National Vulnerability Database (NVD) enriches CVE
entries with severity scores, affected-product configurations, and a weakness
classification~\cite{nvd}. The classification uses the Common Weakness
Enumeration (CWE), a hierarchical catalog of software and hardware weakness
types, such as CWE-22 (path traversal), CWE-79 (cross-site scripting), or
CWE-269 (improper privilege management)~\cite{cwe}. CVE and CWE answer the
questions \emph{what} is vulnerable and \emph{what kind} of flaw it is; they
deliberately do not describe \emph{how} to exploit it. A typical entry
consists of a few sentences of prose and links to an advisory, a patch
commit, and occasionally a proof of concept (PoC).

\subsection{Approaches to Vulnerability Reproduction}

\paragraph{Manual reproduction.}
The traditional approach is manual. An analyst reads the advisory and the
patch, locates and builds the vulnerable version, configures it, and writes
or adapts an exploit, iterating until the exploit works. Mu et
al.~\cite{mu2018reproducibility} quantified the cost of this process: a group
of security analysts spent 3,600 hours attempting to reproduce 368
vulnerabilities, and the information in individual reports was often
insufficient for success. Their survey of 43 practitioners found that
reproduction relies heavily on manual debugging and guesswork. Exploit
databases~\cite{exploitdb} and environment collections such as
Vulhub~\cite{vulhub} reduce part of this effort, but each covers only half of
the problem: one provides exploits without environments, the other
environments without a uniform notion of success or instrumentation.

\paragraph{Automated reproduction: CVE-Genie.}
\genie~\cite{ullah2026cvegenie}, to appear at ACM CCS 2026, automates the
whole process with a multi-agent system built on large language models.
Starting from a CVE identifier, \genie collects the associated resources
(source code at the vulnerable version, advisory, CWE, patch, and PoC when
available), rebuilds the vulnerable project, develops an exploit, and
produces a \emph{verifier}: a script that runs the exploit and checks its
effect. The system is organized in four stages, each made of
developer agents paired with critic agents: a Processor, a Builder, an
Exploiter, and a CTF Verifier. The verifier is designed after Capture-the-Flag
challenges: it prints a fixed flag string only if it observes that the
exploit had its intended effect, so that reproduction success can be checked
mechanically. Evaluated on 841 CVEs published between June 2024 and May 2025,
\genie reproduced 428 of them (about 51\%) across 267 projects, 141 CWEs, and
22 programming languages, at an average cost of \$2.77 per CVE. The authors
release the reproduced CVEs together with the full logs of the agents'
interactions.

\genie shows that reproduction can be automated at scale, but its outputs are
tied to the environment in which the agents worked: a long-lived machine on
which the agents installed packages, compiled code, wrote helper files, and
started services interactively. The reproduced artifacts (the source
snapshot, the exploit, and the verifier) do not capture this state, so they
cannot be rerun elsewhere without reconstructing it. Packaging reproductions
so that they can be rerun anywhere is precisely the problem that \cvex
addresses, and \genie's flag-based verifier is one of the main inspirations
for the success criterion of \cvexthree.

\subsection{Previous Versions of CVEX}

\cvexthree is the fourth iteration of an effort that began at UC Santa
Barbara in 2021. Table~\ref{tab:generations} summarizes the earlier versions.

\begin{table*}[t]
\centering
\caption{Evolution of the CVEX format.}
\label{tab:generations}
\small
\begin{tabular}{@{}lllll@{}}
\toprule
 & \textbf{v1 (2021--22)} & \textbf{v2 (2023--24)} & \textbf{VM framework (2024--25)} & \textbf{v3 (2026)} \\
\midrule
Isolation & Docker Compose & Docker Compose & Vagrant VMs + Ansible & Docker (+ optional services) \\
Roles & target, exploiter, collector & server, client, listener & target VM(s), attacker VM, router VM & target, benign, attacker, recorder \\
Descriptor & \texttt{docker-compose.yml} + \texttt{cvex.json} & \texttt{docker-compose.yml} & \texttt{cvex.yml} + playbooks & \texttt{cvex.yml} + Dockerfiles \\
Success check & file created/read, leak, \texttt{dmesg} & manual inspection & manual inspection & planted flag, verifier, crash witness \\
Observability & pcap & pcap & pcap, TLS (mitmproxy), strace, Procmon & pcap + strace, benign \emph{and} attack \\
Platforms & Linux & Linux & Linux and Windows & Linux \\
Entries & 52 & 6 & 24 & \NumRecords (\NumFunctional functional) \\
\bottomrule
\end{tabular}
\end{table*}

\paragraph{Version 1 (2021--2022).}
The first version defined a \cvex entry as a Docker Compose project with
three services: a \emph{target} running the vulnerable software, an
\emph{exploiter} running the attack, and an optional \emph{collector}, a
\texttt{tcpdump} sidecar sharing the target's network namespace. A JSON
metadata file, validated by a schema, listed the CVEs, the authors, and an
\texttt{exploit\_verification} clause stating how to decide success: a file
created or read in the target, a leaked address, or a kernel log message.
About fifty entries were produced, and a community scoreboard rewarded
contributors with points for CVEs that others had upvoted. Images were
served from a project registry, which has since become unavailable, so most
of these entries can no longer be run as published.

\paragraph{Version 2 (2023--2024).}
A team of undergraduate researchers revisited the format with a simpler
server/client/listener Compose pattern, published prebuilt images on the
GitHub container registry, and catalogued the entries on a website. The
version dropped the metadata file and the formal success check: success was
established by inspecting the captured traffic.

\paragraph{The VM-based framework (2024--2025).}
The following version replaced containers with full virtual machines managed
by Vagrant and provisioned by Ansible, in order to support OS-level and
Windows vulnerabilities. Reusable \emph{blueprints} describe a network of
VMs (e.g., an Ubuntu attacker and an Ubuntu target), and per-CVE
\emph{records} add the playbooks and the attack commands. A router VM
captures all traffic, decrypts TLS with a transparent \texttt{mitmproxy},
and the framework traces system calls with \texttt{strace} on Linux and API
calls with Process Monitor on Windows. Twenty-four records were contributed
through pull requests. The framework is powerful, but heavy: every record
requires multi-gigabyte VM images and a hypervisor, and, like its
predecessors, it has no formal definition of success. Whether an exploit
worked is left to the person inspecting the traces.

\paragraph{Lessons.}
Three lessons from these versions shaped \cvexthree. First, an entry must be
rebuildable from what it contains: entries that depended on a registry or on
prebuilt images disappeared with them. Second, success must be defined by the
record itself and checked by a machine; otherwise a dataset may contain
failed attacks labeled as attacks. Third, the format must be light enough to
run hundreds of records on a single host, which favors containers over VMs
for the large class of vulnerabilities that do not need kernel-level control.
